\documentclass[11pt,a4paper]{ctexart}
\usepackage[a4paper,margin=1.55cm]{geometry}
\usepackage{array}
\usepackage{tabularx}
\usepackage{enumitem}
\usepackage{xcolor}
\usepackage{hyperref}
\usepackage{setspace}
\usepackage{titlesec}
\usepackage{fancyhdr}
\usepackage{tikz}

\definecolor{Accent}{HTML}{1F2937}
\definecolor{TextGray}{HTML}{374151}
\definecolor{LightGray}{HTML}{6B7280}

\hypersetup{colorlinks=true,urlcolor=Accent,linkcolor=Accent}
\pagestyle{fancy}
\fancyhf{}
\renewcommand{\headrulewidth}{0pt}
\renewcommand{\footrulewidth}{0pt}
\cfoot{\textcolor{LightGray}{詹绍基 | 简历}}

\titleformat{\section}{\large\bfseries\color{Accent}}{}{0em}{}
\titlespacing*{\section}{0pt}{9pt}{5pt}
\setlist[itemize]{leftmargin=1.6em,itemsep=2pt,topsep=2pt}
\setlength{\parindent}{0pt}
\linespread{1.14}

\begin{document}

{\LARGE\bfseries\color{Accent} 詹绍基}

\vspace{4pt}
\textcolor{TextGray}{手机：18550954831 \quad 邮箱：\href{mailto:z1309457869@163.com}{z1309457869@163.com} \quad 出生：2000.10 \quad 城市：珠海}

\vspace{4pt}
\textcolor{TextGray}{研究关键词：脉冲神经网络（SNN）、神经可塑性、具身智能、隐私与能效协同优化、可信评测}

\vspace{4pt}
\noindent\textcolor{LightGray}{\rule{\textwidth}{0.5pt}}

\section*{教育背景}
\begin{tabularx}{\textwidth}{>{\bfseries}p{2.2cm}Xr}
本科 & 苏州大学，医学影像学（辅修计算机） & 2019.09--2024.06 \\
二学士 & 北京师范大学，法学（第二学士学位） & 2024.09--2026.07（预计） \\
考试成绩 & 2026考研初试总分329（英语一、数学一、843互联网创新设计），CET-6：595 & \\
\end{tabularx}

\section*{科研与工程能力}
\begin{itemize}
  \item \textbf{SNN建模与训练：} LIF神经元建模、代理梯度训练、STDP/奖励调制可塑性。
  \item \textbf{可信AI评测：} 成员推断攻击（MIA）、DP-SGD扫描、隐私-性能-能效三目标权衡分析。
  \item \textbf{实验工程化：} PyTorch实验流水线、多seed统计、自动化产出CSV/Markdown/HTML/Pareto结果。
  \item \textbf{研究执行：} 可独立完成问题定义、实验设计、结果解释与复现材料整理，可按不同课题组方向快速适配任务。
\end{itemize}

\section*{重点项目经历（独立开发）}
\textbf{1. EmbodiedSNNPrototype：具身闭环脉冲神经网络原型}
\begin{itemize}
  \item 搭建 environment $\rightarrow$ sensory encoding $\rightarrow$ LIF brain $\rightarrow$ motor readout $\rightarrow$ environment 的闭环系统。
  \item 实现 structured / structured\_plastic / rewired / random\_sparse 多连接模式与参数网格搜索。
  \item 最优structured\_plastic配置在综合目标上显著领先（objective\_mean=2.0327）。
  \item 同时达到较高行为指标：food\_eaten\_mean=3.24，near\_food\_ratio\_mean=0.5367。
\end{itemize}

\textbf{2. Neuro-Symbiosis：SNN-Transformer脑机解码与三目标协同优化}
\begin{itemize}
  \item 构建SNN前端 + Transformer时序解码后端，面向BCI任务联合优化效用、隐私与能效。
  \item 打通训练、baseline对比、隐私评估、能耗评估、DP-SGD sweep、多seed鲁棒性评测全链路。
  \item 形成可证伪问题：spike rate与有效token长度如何共同影响总能耗及隐私-性能权衡。
\end{itemize}

\textbf{3. MedSparseSNN：医疗影像中的稀疏SNN隐私-效率研究}
\begin{itemize}
  \item 在DermaMNIST与PathMNIST对比SNN / DenseSNN / ANN的准确率、隐私风险与效率。
  \item PathMNIST：SNN test\_acc=82.33$\pm$0.31（ANN 85.12$\pm$0.40），理论有效MAC节省84.18\%$\pm$0.52\%。
  \item DermaMNIST：SNN test\_acc=69.93$\pm$0.20，理论有效MAC节省90.74\%$\pm$0.23\%。
\end{itemize}

\textbf{4. PrivEnergyBench：隐私-能效-效用三目标评测工具箱}
\begin{itemize}
  \item 将训练、MIA、时延、能耗代理整合为统一benchmark工具链，支持linear/MLP/CNN1D/Transformer。
  \item 默认3-seed结果：MLP最高val\_acc=0.9208；Linear最低能耗0.0164 mJ/sample。
\end{itemize}

\section*{补充经历与证明}
\begin{itemize}
  \item 上海交通大学医学院苏州九龙医院临床实习（2023.06--2024.06）：完成多科室轮转，出科考试平均93分，完成影像报告300+份，综合合格率96\%。
  \item 证明材料可提供：本科毕业证/学位证/成绩单、学信网学籍学历验证报告、学位在线验证报告、考研初试成绩证明。
  \item GitHub：\href{https://github.com/ZhanMous}{https://github.com/ZhanMous}
\end{itemize}

\end{document}
